\section{The sharp boundary for a fixed native refinement}\label{sec:capacity}
\subsection{Quadratic type and its Gauss sign}
Let $(V,B)$ be a nondegenerate alternating space over $\F_2$ of dimension $2s$, and let $q$ be a quadratic refinement. In a symplectic basis $(e_j,f_j)$,
\[
 q(x_1,y_1,\ldots,x_s,y_s)
 =\sum_{j=1}^s\bigl(x_jy_j+q(e_j)x_j+q(f_j)y_j\bigr).
\]
Its Gauss sum factors into plane sums:
\[
 G(q)=\sum_{v\in V}(-1)^{q(v)}
 =2^s(-1)^{\sum_j q(e_j)q(f_j)}.
 \tag{27.1}\label{eq:gauss}
\]
Define $\Arf(q)\in\F_2$ by the sign in \eqref{eq:gauss}. The sum is intrinsic, so this definition is independent of the symplectic basis.

\begin{lemma}[Classification in the required finite setting]\label{lem:arfclass}
Two nondegenerate quadratic spaces over $\F_2$ are isometric if and only if they have the same dimension and Arf invariant. Both Arf types occur in every positive even dimension.
\end{lemma}
\begin{proof}
The coordinate expression decomposes a space into nondegenerate planes. On a plane, the three refinements with Arf invariant zero are isometric to $H(x,y)=xy$, while the unique remaining type is $A(x,y)=xy+x+y$. This can be checked on its three nonzero vectors.

Two anisotropic planes can be replaced by two hyperbolic planes. In $A\perp A$, the vector $u=e_1+e_2$ is nonzero and isotropic. Choose $v$ with $B(u,v)=1$. Replacing $v$ by $v+q(v)u$ makes it isotropic too, so $u,v$ span a hyperbolic plane. Its nondegenerate orthogonal complement is a plane whose Gauss sign is positive by multiplicativity of \eqref{eq:gauss}. It too is hyperbolic. Pairing all anisotropic planes leaves either only hyperbolic planes or one anisotropic plane. The parity of the remaining anisotropic planes is exactly the Arf invariant. These normal forms prove classification and existence.
\end{proof}

In dimension six, \eqref{eq:gauss} gives $G(q)=+8$ or $-8$. The numbers of zero values are respectively
\[
 \#q^{-1}(0)=\frac{64+8}{2}=36,
 \qquad \#q^{-1}(0)=\frac{64-8}{2}=28.
 \tag{27.2}\label{eq:rank8zeros}
\]
These counts include the zero vector. The new suite verifies both refinements on all 64 vectors. They supply a concrete signature of the distinction that dimension alone misses.

\subsection{One spare plane is enough only when the embedding is free}
\begin{proposition}[Fixed-refinement capacity criterion]\label{prop:capacity}
Let $(V,q_V)$ and $(W,q_W)$ be nondegenerate quadratic spaces of dimensions $2s$ and $2R$. When the embedding may be freely chosen, a quadratic isometric embedding $V\to W$ exists exactly when
\[
 R\ge s,
 \qquad R=s\ \Longrightarrow\ \Arf(q_V)=\Arf(q_W).
 \tag{27.3}\label{eq:capcrit}
\]
If $R>s$, either source Arf type embeds.
\end{proposition}
\begin{proof}
An embedding requires the dimension inequality. At equal dimension it is an isometry, so the Arf invariants must agree. If $R>s$, choose a nondegenerate space $U$ of dimension $2(R-s)$ with
\[
 \Arf(q_U)=\Arf(q_W)+\Arf(q_V).
\]
Both choices are available by Lemma~\ref{lem:arfclass}. The space $V\perp U$ has the same dimension and Arf invariant as $W$, hence is isometric to $W$. Restrict that isometry to $V$. The equal-dimension sufficiency follows from the same classification.
\end{proof}

After the quadratic embedding is chosen, the preimage of its image in an actual $\mu_2$-central extension of $W$ has the source refinement. Proposition~\ref{prop:ext} supplies the extension isomorphism. Exact generating lifts can be matched when their squares and commutators match. Additional prescribed section values retain the stronger product condition from Section~\ref{sec:extensions}.

The spare-plane argument does not permit arbitrary changes to fixed depth placements, transport maps, or stabilizer actions. It proves existence of an otherwise free quadratic embedding. These qualifications are important even when the host has a large dimensional surplus.

\subsection{The exact rank-eight consequence}
At $d=3$ in \eqref{eq:dyadic-D}, the full dyadic quotient and the three-plane source both have dimension six. Suppose the source planes have Arf invariants $\epsilon_1,\epsilon_2,\epsilon_3$, and a particular actual host refinement has Arf invariant $\epsilon_W$. Then the necessary and sufficient test for a \emph{free} quadratic embedding is
\[
 \epsilon_W=\epsilon_1+\epsilon_2+\epsilon_3\quad\hbox{in }\F_2.
 \tag{27.4}\label{eq:rank8arf}
\]
There is no orthogonal complement in which to absorb a mismatch. For $d\ge4$, the full profile has a positive-dimensional complement, so the finite Arf obstruction disappears when the embedding is free and the actual host refinement is supplied.

This does not compute $\epsilon_W$ for a particular arithmetic type. A Heisenberg commutator quotient alone supplies an alternating form. It does not canonically select a finite two-element-center extension or its square refinement. Determining an independently required native refinement and its lift phases for the chosen candidate remains an arithmetic obligation.

Conversely, under the transported scalar-pushout contract of Section~\ref{sec:extensions}, the host refinement is part of the transported output rather than independently prescribed. Every source marking can then be retained in the same scalar extension once the symplectic capacity and recovery hypotheses hold. There is no contradiction between this positive statement and \eqref{eq:rank8arf}. They quantify over different data.

\source{The capacity criterion was recorded in [Q16]. Its finite classification input and proof are made explicit here. The rank-eight deduction combines that criterion with [D167]. The block does not assign a canonical native Arf invariant to an actual arithmetic candidate that has not supplied one.}
